\documentclass[11pt,a4paper]{article}

\usepackage[a4paper,margin=24mm]{geometry}
\usepackage{fontspec}
\usepackage[english]{babel}
\usepackage{hyperref}
\usepackage{longtable}
\usepackage{array}
\usepackage{xcolor}
\usepackage{graphicx}
\usepackage{float}

\setmainfont{Helvetica Neue}
\emergencystretch=3em
\sloppy
\hypersetup{
  colorlinks=true,
  linkcolor=blue,
  urlcolor=blue,
  pdftitle={musiccloud Resolver Flow},
  pdfauthor={musiccloud}
}

\newcommand{\code}[1]{\texttt{#1}}
\newcommand{\versionstamp}{Version DOC\_VERSION, DOC\_DATE}

% Inline screenshot with a caption. Bounded in width and height so tall
% captures never overflow a page; aspect ratio is preserved.
\newcommand{\shot}[2]{%
  \begin{center}
  \includegraphics[width=0.92\linewidth,height=0.62\textheight,keepaspectratio]{#1}\\[3pt]
  {\small\itshape #2}
  \end{center}}

\title{musiccloud Resolver Flow}
\author{musiccloud}
\date{\versionstamp}

\begin{document}
\maketitle
\tableofcontents
\newpage

\section{Overview}

The resolver accepts a link or search query and turns it into a unified musiccloud response. The response is either a resolved track, album, or artist, a candidate list for an ambiguous track search, a genre-discovery result, or an error with a stable error code.

In the musiccloud web app, every query starts at a single input. The screenshot below shows that entry point: it accepts a streaming-service URL, free text, a structured \code{title:}/\code{artist:} query, or a \code{genre:} discovery query.

\shot{screenshots/01-hero-input.png}{The single input that accepts every query shape.}

This document describes the public query shapes from a user's perspective. Implementation details such as enabled plugins, admin toggles, or internal cache tables are not required to use the API. The screenshots illustrate how each query shape appears in the web app; the JSON blocks show the matching API request and response.

\section{Endpoints}

\begin{longtable}{>{\raggedright\arraybackslash}p{0.24\linewidth}>{\raggedright\arraybackslash}p{0.68\linewidth}}
\textbf{Endpoint} & \textbf{Use} \\
\hline
\code{POST /api/v1/resolve} & Full resolve endpoint for API clients with an \code{X-API-Key} or bearer token. Supports URLs, free text, genre discovery, structured search, and the \code{selectedCandidate} follow-up. \\
\code{GET /api/v1/resolve} & Unauthenticated scripting endpoint for one-shot track resolves. Supports URLs, free text, and structured search. Ambiguous searches return \code{400} because the follow-up state is not carried. Genre discovery is not supported here. \\
\end{longtable}

\section{Query Shapes}

\subsection{Streaming-Service URL}

Use this shape when you already have a track, album, or artist link.

\begingroup\small\begin{verbatim}
https://open.spotify.com/track/2WfaOiMkCvy7F5fcp2zZ8L
https://music.apple.com/de/album/example/123456789
open.spotify.com/track/2WfaOiMkCvy7F5fcp2zZ8L
https://link.deezer.com/s/abc123
\end{verbatim}\endgroup

The resolver detects the service, reads the source metadata, and searches for matching links on other platforms. Schemeless host forms (without \code{https://}) are accepted for allow-listed music hosts. Short links (for example \code{link.deezer.com}) are expanded first. Tracking parameters are removed before persistence and before the response is returned.

A request and its (trimmed) response look like this. The \code{links} array carries one entry per platform the track was matched on, each with a \code{confidence} and the \code{matchMethod} that produced it (\code{isrc}, \code{upc}, \code{search}, \code{cache}, or \code{isrc-inference}):

\begingroup\small\begin{verbatim}
GET /api/v1/resolve
  ?query=https://open.spotify.com/track/2WfaOiMkCvy7F5fcp2zZ8L

200 OK
{
  "id": "I74PsoH_M-oDufeUQW3WT",
  "shortUrl": "https://musiccloud.io/THmwY",
  "track": {
    "title": "Take on Me",
    "artists": ["a-ha"],
    "albumName": "Hunting High and Low",
    "artworkUrl": "https://i.scdn.co/image/ab67616d0000b273e8dd...",
    "durationMs": 225280,
    "isrc": "USWB19901214",
    "releaseDate": "1985-06-01",
    "isExplicit": false,
    "previewUrl": "https://cdnt-preview.dzcdn.net/api/1/.../...mp3"
  },
  "links": [
    { "service": "spotify", "displayName": "Spotify",
      "url": "https://open.spotify.com/track/2WfaOiMkCvy7F5fcp2zZ8L",
      "confidence": 1, "matchMethod": "isrc" },
    { "service": "apple-music", "displayName": "Apple Music",
      "url": "https://music.apple.com/us/album/take-on-me/747796769",
      "confidence": 1, "matchMethod": "isrc" }
  ]
}
\end{verbatim}\endgroup

In the web app, the same resolve renders as a result card with the cover, a preview player, a share action, and the cross-service ``Listen on'' links:

\shot{screenshots/02-resolved-result.png}{A resolved track with cross-service links and artist info.}

On \code{POST}, album and artist URLs resolve as well and return \code{type: "album"} or \code{type: "artist"} instead of a \code{track}. The \code{GET} endpoint resolves track URLs only.

\subsection{Free-Text Query}

Free text is any input that is not a URL and does not start with a reserved query prefix.

\begingroup\small\begin{verbatim}
bohemian rhapsody queen
radiohead karma police ok computer
daft punk one more time
\end{verbatim}\endgroup

This shape is useful when the user only roughly knows what they want. The whole string is used as both the title and the artist term for each adapter; there is no client-side split. A high-confidence match resolves directly to a track response. If there are multiple plausible matches, \code{POST /api/v1/resolve} returns a candidate list (see below), while \code{GET} returns \code{400} because it cannot carry the follow-up.

\begingroup\small\begin{verbatim}
POST /api/v1/resolve
{ "query": "daft punk one more time" }

200 OK
{
  "status": "disambiguation",
  "candidates": [
    { "id": "spotify:0DiWol3AO6WpXZgp0goxAV",
      "title": "One More Time", "artists": ["Daft Punk"],
      "albumName": "Discovery",
      "artworkUrl": "https://i.scdn.co/image/ab67616d0000b273..." },
    { "id": "spotify:1hLNqXjr5371Rd3llz5Gof",
      "title": "One More Time (Short Radio Edit)", "artists": ["Daft Punk"],
      "albumName": "One More Time" }
  ]
}
\end{verbatim}\endgroup

The web app shows that candidate list as a ``Did you mean?'' panel with up to eight rows per page:

\shot{screenshots/03-disambiguation.png}{An ambiguous free-text search resolves to a candidate list.}

\subsection{Genre-Discovery Query}

Genre discovery always starts with \code{genre:}. It is meant for browsing, not for finding one specific track, and is supported on \code{POST} only.

\begingroup\small\begin{verbatim}
genre: ?
genre: jazz
genre: jazz, tracks: 20, vibe: mixed
genre: jazz|r&b, count: 5
\end{verbatim}\endgroup

\code{genre: ?} returns a browse grid of popular genres. Each tile carries a \code{name} that you feed back into a \code{genre: <name>} query:

\begingroup\small\begin{verbatim}
POST /api/v1/resolve
{ "query": "genre: ?" }

200 OK
{
  "status": "genre-browse",
  "genres": [
    { "name": "jazz", "displayName": "Jazz",
      "artworkUrl": "/api/genre-artwork/jazz?v=5",
      "accentColor": "#3b2f63" },
    { "name": "synthpop", "displayName": "Synthpop",
      "artworkUrl": "/api/genre-artwork/synthpop?v=5" }
  ]
}
\end{verbatim}\endgroup

\shot{screenshots/04-genre-browse.png}{The \texttt{genre: ?} browse grid of popular genres.}

\code{genre: <name>} returns up to three parallel result lists for tracks, albums, and artists. Each list may be \code{null} when that type was not requested. The optional modifiers are \code{tracks}, \code{albums}, \code{artists}, \code{count}, and \code{vibe} (\code{hot} or \code{mixed}). Track and album rows carry a \code{webUrl} that you feed into a follow-up resolve to produce a result card.

\begingroup\small\begin{verbatim}
POST /api/v1/resolve
{ "query": "genre: jazz, tracks: 10, artists: 3" }

200 OK
{
  "status": "genre-search",
  "query": {
    "genres": ["jazz"], "vibe": "hot",
    "tracks": 10, "albums": null, "artists": 3
  },
  "results": {
    "tracks": [
      { "id": "deezer:1109731", "title": "Take Five",
        "artists": ["The Dave Brubeck Quartet"],
        "albumName": "Time Out", "durationMs": 324000,
        "webUrl": "https://www.deezer.com/track/1109731" }
    ],
    "albums": null,
    "artists": [
      { "id": "deezer:1716", "name": "Miles Davis",
        "webUrl": "https://www.deezer.com/artist/1716" }
    ]
  },
  "warnings": []
}
\end{verbatim}\endgroup

\shot{screenshots/05-genre-results.png}{A \texttt{genre:} search renders as three result columns.}

\section{Structured Search}

Structured search starts with \code{title:}, \code{artist:}, or \code{album:}. It is meant for concrete track searches where individual fields are known.

\subsection{Purpose}

Structured search is more precise than free text because title, artist, and optional album can be passed to adapters separately. Services with field operators, such as Spotify, Deezer, or MusicBrainz, can search more narrowly. Other services receive the fields as combined search terms.

\subsection{Syntax}

\begingroup\small\begin{verbatim}
title: The Killing Moon, artist: Echo & The Bunnymen
artist: Radiohead
title: Karma Police, artist: Radiohead, album: OK Computer, count: 5
\end{verbatim}\endgroup

\begin{longtable}{>{\raggedright\arraybackslash}p{0.18\linewidth}>{\raggedright\arraybackslash}p{0.72\linewidth}}
\textbf{Field} & \textbf{Meaning} \\
\hline
\code{title} & Track title. At least one of \code{title} or \code{artist} is required. \\
\code{artist} & Artist name. At least one of \code{title} or \code{artist} is required. \\
\code{album} & Optional album title for refinement. \code{album} alone is rejected. \\
\code{count} & Optional integer from 1 to 10. Caps the candidate list when the search is ambiguous. \\
\end{longtable}

Field names are case-insensitive. Values preserve their original casing. Fields are separated by commas. A missing comma before the next known field is tolerated, for example \code{title: foo artist: bar}. Commas inside values are not escaped and act as field separators.

\subsection{Validation}

Empty values, duplicate fields, unknown fields, and \code{count} outside 1 to 10 return \code{400}. Genre modifiers such as \code{tracks}, \code{albums}, \code{artists}, and \code{vibe} are intentionally rejected in structured search and belong in \code{genre:} queries.

\subsection{Response}

With high confidence, the resolver returns a track response directly, identical in shape to the URL resolve above. With multiple plausible matches, \code{POST /api/v1/resolve} returns a candidate list. \code{count} caps this list in addition to the server-side maximum. \code{GET /api/v1/resolve} cannot continue a candidate selection and returns \code{400} when the query is ambiguous:

\begingroup\small\begin{verbatim}
GET /api/v1/resolve?query=artist:%20Radiohead

400 Bad Request
{
  "error": "MC-URL-0003",
  "message": "Structured query was ambiguous; use POST endpoint
              for disambiguation. (MC-URL-0003)"
}
\end{verbatim}\endgroup

\section{selectedCandidate Follow-Up}

When a search is ambiguous, \code{POST /api/v1/resolve} returns a response with \code{status: "disambiguation"} and candidates. Each candidate contains an opaque \code{id}. Send that \code{id} back unchanged as \code{selectedCandidate}, and the resolver returns the final track response:

\begingroup\small\begin{verbatim}
POST /api/v1/resolve
{ "selectedCandidate": "spotify:0DiWol3AO6WpXZgp0goxAV" }

200 OK
{
  "id": "Xa3kPq7r-T0mfZ2Lc9D1S",
  "shortUrl": "https://musiccloud.io/9Kdp2",
  "track": {
    "title": "One More Time", "artists": ["Daft Punk"],
    "albumName": "Discovery", "isrc": "GBDUW0000059"
  },
  "links": [
    { "service": "spotify", "displayName": "Spotify",
      "url": "https://open.spotify.com/track/0DiWol3AO6WpXZgp0goxAV",
      "confidence": 1, "matchMethod": "isrc" }
  ]
}
\end{verbatim}\endgroup

Clients should not construct the ID themselves. The format is internal and may vary by service.

\section{Error Cases}

Errors carry a stable, machine-readable \code{error} code and a human-readable \code{message} that ends with the same code in parentheses:

\begingroup\small\begin{verbatim}
404 Not Found
{
  "error": "MC-RES-0001",
  "message": "This track doesn't seem to be available anymore
              on the source service. (MC-RES-0001)"
}
\end{verbatim}\endgroup

\begin{longtable}{>{\raggedright\arraybackslash}p{0.18\linewidth}>{\raggedright\arraybackslash}p{0.72\linewidth}}
\textbf{Status} & \textbf{Typical Cause} \\
\hline
\code{400} & Malformed query, unknown field, empty structured-search values, ambiguous GET search, or invalid \code{selectedCandidate}. \\
\code{401} & Missing or invalid credentials on protected endpoints. \\
\code{404} & The query is valid, but no matching track, album, or artist was found. \\
\code{408} & An upstream service did not answer in time. \\
\code{429} & Rate limit exceeded. \\
\code{503} & Required upstream service is temporarily unavailable. \\
\end{longtable}

\section{Practical Examples}

\begingroup\small\begin{verbatim}
# Resolve a Spotify track URL (one-shot, unauthenticated)
GET /api/v1/resolve
  ?query=https://open.spotify.com/track/2WfaOiMkCvy7F5fcp2zZ8L

# Structured search, capped to 3 candidates
POST /api/v1/resolve
{ "query": "title: Karma Police, artist: Radiohead, count: 3" }

# Genre discovery: 10 tracks and 5 albums of jazz
POST /api/v1/resolve
{ "query": "genre: jazz, tracks: 10, albums: 5" }

# Continue a disambiguation by selecting a candidate
POST /api/v1/resolve
{ "selectedCandidate": "deezer:3135556" }

# Free text via the scripting endpoint (URL-encoded)
GET /api/v1/resolve
  ?query=title:%20Bohemian%20Rhapsody,%20artist:%20Queen
\end{verbatim}\endgroup

\end{document}
